As large pre-trained language models become increasingly commonplace, their flaws have likewise become increasingly glaring.
    %%CF.8.30: this seems like more of an opinion -- the flaws of older models were also super glaring because they didn't work! Could rephrase to something like "Large pre-trained language models have delivered impressive results on a range of downstream tasks, but often have bugs or errors that are discovered after the fact: their outputs may become outdated over time or be wrong in the first place for some inputs."
    While powerful, large pre-trained models can fail in surprising and even dangerous ways; their outputs become outdated, reproduce harmful stereotypes or biases present in training data,
    %%CF.8.30: I think we shouldn't talk about stereotypes or biases unless we explicitly include results on this in the paper. Also, if you talk about stereotypes and biases in the abstract, then paper matching will highly rank people who work on that.
    or are simply wrong.
    %%CF.8.30: superficially convincing is also an opinion, and not always actually true. It's also not particularly important/relevant to this paper. EM: Done
    Ideally, a developer of such a model might also release tools to enable end-users to make targeted corrections
    %%CF.8.30: the developer themself may want to make edits too?
    to correct undesirable outputs while leaving the model otherwise intact. However, the distributed, black-box nature of the representations learned by large neural networks makes such targeted edits difficult. While existing work has proposed several approaches to enable targeted model edits, the most effective approaches rely on a bi-level optimization that incurs very high computational costs for large models. 
    %%CF.8.30: I wouldn't bring up ENNs in the abstract because 99% of your readers won't be familiar with them -- they won't be the obvious solution that the reader thinks of to this problem. I would instead say that fine-tuning or re-training the model from scratch each time a mistake is found quickly becomes inordinately expensive for large models.
    Further, we find that applying such bi-level approaches to large pre-trained models in natural language settings often induces a drop in performance (e.g., increase in perplexity) even \textit{before} an edit has been applied. Given these algorithmic problems, we propose \name, an algorithm leveraging the low-rank structure of the gradients in neural networks to perform efficient, accurate edits for any pre-trained model, irrespective of initialization.
    %%CF.8.30: I think that some of these details aren't going to make sense to the reader without context. For things like that, I would omit them and be more vague rather than be confusing. For example: "Motivated by this problem, we propose to train an editing network that can produce fast local changes to a model based on specified input, output change."
    Our experiments show that \sname~is as or more effective at editing large pre-trained transformer language models than prior approaches to model editing across a range of problems, and can scale to models with over a billion parameters without modification.